%
% 9-fazit.tex -- Komplexe Eigenwerte als Entwurfswerkzeug
%
% (c) 2026 Stephanie Märklin, OST Ostschweizer Fachhochschule
%
% !TEX root = ../../buch.tex
% !TEX encoding = UTF-8
%
% Rolle im Paper: Fazit
% Roter Faden: Das dynamische Verhalten steckt in wenigen komplexen Zahlen,
% anwendbar im Entwurf statt im Flugversuch.
% Überleitung aus dem vorangehenden Abschnitt: Was bleibt als Ergebnis?
%
\section{Komplexe Eigenwerte als Entwurfswerkzeug
\label{aircraft:section:fazit}}
\kopfrechts{Entwurfswerkzeug}
Die Leitfrage lautet, warum wir die Stabilität eines Flugzeugs berechnen
müssen und wie wir es können.
Berechnen müssen wir die Stabilität, weil die Beobachtung im Flug erst möglich
ist, wenn das Flugzeug existiert, und weil die Beobachtung kein Mass für die
Stabilität liefert (Abschnitt~\ref{aircraft:section:3-mathematische-stabilitaet}).
Berechnen können wir die Stabilität, weil die Seitenbewegung für kleine
Störungen dem linearen Bewegungsgesetz \eqref{aircraft:eqn:bewegungsgesetz}
folgt.
Stabilität bedeutet in diesem Modell, dass jede Lösung abklingt
(Abschnitt~\ref{aircraft:section:3-mathematische-stabilitaet}).

Der Exponentialansatz \eqref{aircraft:eqn:ansatz} führt das Bewegungsgesetz
auf die Eigenwertgleichung \eqref{aircraft:eqn:eigenwertgleichung} zurück.
Die Eigenwerte sind die Nullstellen der charakteristischen Gleichung
\eqref{aircraft:eqn:charakteristisch}.
Der Realteil eines Eigenwerts entscheidet nach
Abschnitt~\ref{aircraft:section:5-interpretation}, ob die zugehörige Bewegung
abklingt, der Imaginärteil, ob sie schwingt.
Halbwertszeit \eqref{aircraft:eqn:halbwertszeit} und Periode übersetzen die
Eigenwerte in Zeiten, die der Pilot im Flug wahrnimmt.
Das dynamische Verhalten der Seitenbewegung ist damit durch vier komplexe
Zahlen bestimmt.

Die Derivative in der Systemmatrix liegen aus der Geometrie vor, bevor das
Flugzeug fliegt (Abschnitt~\ref{aircraft:subsection:systemmatrix}).
Die Eigenwerte lassen sich deshalb im Entwurf berechnen und mit den
Anforderungen an die Flugeigenschaften vergleichen.
Die Anforderungen, die Nelson \cite{aircraft:nelson} in Tabelle~5.4
wiedergibt, verlangen für Flugzeuge mittleren Gewichts eine Verdopplungszeit
der Spiralbewegung von mindestens $20$\,s.
Der Business Jet aus Abschnitt~\ref{aircraft:section:6-beispiel} erfüllt diese
Anforderung mit $78$\,s, obwohl seine Spiralbewegung instabil ist.
Genügt ein Entwurf den Anforderungen nicht, so zeigt
Abschnitt~\ref{aircraft:section:fluegel}, welche Derivative zu verändern sind,
und Abschnitt~\ref{aircraft:section:grenzen}, wie ein Gierdämpfer die Dämpfung
künstlich erhöht.
Die Frage, ob ein Flugzeug stabil ist, beantworten wir damit, bevor das
Flugzeug gebaut ist.
